\chapter{基金设立：载体选择与基本条款}

\section{有限合伙：PE 的主流选择}

在多数法域，\textbf{有限合伙}因其\textbf{税收穿透}（在符合条件下由合伙人层面纳税而非实体层面双重征税）、GP 无限责任与 LP 有限责任的清晰划分，成为私募股权基金的首选载体。基金期限、投资期与延长期、GP 移除与关键人条款等，均在 LPA 中锁定。设立成本包括注册费用、首次法律起草与监管登记（若适用）；跨境 LP 时还需考虑反洗钱（KYC）与制裁筛查流程。

\section{公司制与契约型：何时出现}

\textbf{公司制}基金在部分辖区便于某些类型投资者（如养老基金子账户）持有，或用于上市载体（如部分投资信托结构）；劣势常在于\textbf{双重征税}或治理灵活性不足。\textbf{契约型基金}在部分亚洲市场与资管产品中常见，依赖托管人与管理人合同网络；与经典离岸/在岸有限合伙相比，权利义务映射方式不同，需在律师协助下逐条对照 LPA 等价物。

\section{LPA 中的经济条款框架}

\textbf{管理费}通常按认缴或实缴或 NAV 的一定比例计提，投资期后可能调降或切换基数。\textbf{业绩报酬}在超过 hurdle 后对利润分成；需明确是\textbf{whole-of-fund}还是 deal-by-deal、是否存在 clawback（回拨）机制。\textbf{分配瀑布}决定税、费、本金与收益在 LP 与 GP 之间的先后顺序，直接影响 IRR 与 DPI 的时间形态。Side letter 可能对大型 LP 给予最惠国（MFN）、共同投资权或额外信息权，管理公司需建立台账避免承诺冲突。

\section{治理与管理条款}

投资委员会组成、关联交易审批、个人跟投（co-investment）规则、借款与担保限额、关键人事件与 GP 替换程序，构成管理条款的骨架。从零到一阶段，常见错误是\textbf{过度照搬大型基金模板}导致小团队无法合规执行；应在「可运营」与「LP 可接受」之间折中，并预留随基金规模扩大而修订的修正案机制（须符合适用法律）。
